\section{Introduction}
\label{sec:intro}
Cybersickness---nausea, disorientation and oculomotor strain caused by the sensory conflict of virtual reality (VR)---limits the use of head-mounted displays in rehabilitation, training and assessment~\cite{caserman2021cybersickness}. It can also be induced deliberately by manipulating frame pacing, camera motion or field of view~\cite{tseng2022dark,odeleye2021framerate}, and attacks on a detector's inputs can suppress a mitigation~\cite{kundu2025securing}. A detector that recognizes cybersickness during immersion is therefore a safety function with unequal error costs: a missed event leaves the user exposed, whereas a false alarm causes a brief interruption.

Detectors have moved from post-exposure questionnaires to deep models of physiological, eye-tracking and gameplay data~\cite{dennison2016physio,islam2020automatic,islam2021fusion,liu2024cnneca,yalcin2024automatic,jeong2024precyse}. To encode unequal costs, several authors cast classification as a sequential decision problem for reinforcement learning (RL) and trained deep Q-networks (DQNs) with class-dependent rewards~\cite{lin2020imbalanced,yang2024drlimbalanced}. The enhanced Double Deep Q-Network (E-DDQN) follows this line: a dueling Double DQN~\cite{vanhasselt2016double,wang2016dueling} whose reward penalizes a missed event five times more than a false alarm, trained with class-balanced replay and soft target updates. Unlike~\cite{lin2020imbalanced,yang2024drlimbalanced}, it does not end an episode when an event is missed.

Two questions must be answered before such a detector can be credited with a safety benefit. \emph{Evaluation:} gameplay telemetry is logged about once per second and self-reported discomfort changes rarely within a session, so a split that places seconds of one session in both training and test data rewards recognition of the session~\cite{kapoor2023leakage}; published results on the dataset used here come from a 70/30 split for which no session separation is reported~\cite{kundu2022truvr}, and an earlier, unpublished version of this study reported \pct{96.88} accuracy for E-DDQN after oversampling before splitting. \emph{Formulation:} what an asymmetric reward adds to cost-sensitive learning~\cite{elkan2001cost} is rarely made explicit, and RL detectors are seldom compared with cost-weighted models tuned and thresholded under the same rules.

On the public UFFCSData dataset~\cite{porcino2019uffcsdata} we ask:
\begin{itemize}
\item[RQ1] What does the asymmetric reward add, formally, to cost-sensitive classification?
\item[RQ2] Under leakage-controlled, session-grouped evaluation, does E-DDQN outperform a standard DQN, cost-weighted supervised models and an exposure-time baseline, and does any model meet a safety-prioritized operating target?
\item[RQ3] As a secondary sanity check, how sensitive are the detectors to bounded manipulation of the telemetry channels? The dataset contains no manipulated sessions, so this is a feature-space stress test, not an evaluation of attacks.
\end{itemize}

The contributions are:
\begin{enumerate}
\item \emph{Formulation.} Applying known results~\cite{elkan2001cost,dudik2011efficient}, we show that on logged data the alarm cannot change future observations, so detection is a contextual bandit and the 5:1 reward induces exactly the Bayes rule $P(y=1\mid s)\ge1/6$ of 5:1 cost-sensitive classification; class-balanced replay shifts the operating point actually learned (Section~\ref{sec:formulation}).
\item \emph{A matched, leakage-controlled benchmark.} E-DDQN, a standard DQN, a contextual-bandit variant (E-CB) and four cost-weighted supervised models (a cost-sensitive multilayer perceptron, CS-MLP; logistic regression, LogReg; histogram gradient boosting, HGB; and a random forest, RF) are compared under one nested, session-linkage-grouped protocol fixed before testing, with the same search procedure, number of configurations (not coverage: the search spaces differ greatly in size), checkpoint and threshold rules (Figs.~\ref{fig:workflow}--\ref{fig:protocol}); two exposure-time baselines were added post hoc. Row-level splits of the same data give an area under the receiver operating characteristic curve (ROC-AUC) of \LeakRowRocAuc--\LeakOldRocAuc, against \LeakGrpRocAuc{} under session-grouped splits (Section~\ref{sec:leak}).
\item \emph{Findings.} E-DDQN did not have a significantly lower 5:1 cost than the standard DQN with its primary or two further, pre-specified search draws. Its seed-ensemble cost was higher than that of a cost-weighted random forest, significantly under the primary group bootstrap but not under an exact sign-flip sensitivity test, so this difference is not firmly established. No primary model ranked seconds significantly better than elapsed time alone, and no model combined high recall with a controlled false-alarm rate (Section~\ref{sec:results}).
\item \emph{Diagnosis.} Across three search draws E-DDQN's within-fold ROC-AUC ranged from \SDEDDQNWfMin{} to \SDEDDQNWfMax, which suggests that sensitivity to the search draw contributes substantially to its weak primary result; the draw, feature sets and other hyper-parameters varied together, so no single cause is isolated (Section~\ref{sec:why}).
\item \emph{Reproducibility.} Code, the fixed protocol, per-sample predictions and executable checks of the splits and the replay buffer are released.
\end{enumerate}
Section~\ref{sec:related} reviews related work, Section~\ref{sec:method} presents the methods and setup, Section~\ref{sec:results} the results and Section~\ref{sec:discussion} the discussion; Section~\ref{sec:conclusion} concludes.
